\begin{recipe}{Eierkoeken}
\kicker[Bakken,Zoet,Brabants]{Bakken · zoet · Brabants}
\index[register]{Eierkoeken}
\index[register]{Eieren}
\index[register]{Citroen}

\meta{25 minuten \dvd 7 stuks \dvd Oven 200\,°C}

\lettrine{B}{ij ons} gaan eierkoeken in het weekend weleens mee bij de lunch,
of als tussendoortje onderweg. Het citroen geeft ze een frisse smaak, en het
luchtig kloppen van de eieren zorgt voor de lichte structuur.

\blockrule{Ingrediënten}
\begin{ingredients}
  \ing{2}{eieren}\index[register]{Eieren}
  \ing{100 g}{suiker}
  \ing{120 g}{bloem}
  \ing{8 g}{bakpoeder (½ zakje)}
  \ingb{rasp van 1 citroen}\index[register]{Citroen}
  \ingb{stuifbloem, voor de bakplaat}
  \ingb{snuf zout}
\end{ingredients}

\blockrule{Bereiding}
\tip{Verwarm de bakplaat voor in de oven. Een koude plaat geeft platte koeken, een hete
plaat zorgt dat ze meteen beginnen te rijzen.}
\begin{steps}
  \item Verwarm de oven voor op 200 graden (boven- en onderwarmte) en zet de bakplaat alvast in de oven.
  \item Klop de eieren met de suiker minstens 5 minuten goed luchtig, tot de massa
        echt bleek en dik is en van de garde in een lint valt. Er zit weinig
        bakpoeder in dit beslag, dus dat luchtig geklopte ei is hier de hele
        structuur.
  \item Voeg de citroenrasp, bloem, bakpoeder en een snuf zout toe. Spatel alles voorzichtig
        met de hand door het beslag, niet te veel roeren.
  \item Haal de hete bakplaat uit de oven en bestrooi royaal met stuifbloem.
  \item Schep met twee eetlepels flinke hoopjes beslag op de plaat. Laat voldoende ruimte
        tussen de hoopjes, ze lopen nog iets uit en rijzen.
  \item Bak de eierkoeken 10--12 minuten tot ze goudgeel zijn en licht veren als je erop drukt.
\end{steps}
\end{recipe}
